\section{Related Work}
\label{sec:related}

\para{Validation of harm judges and classifiers.}
\jbb \citep{chao2024jailbreakbench}, HarmBench \citep{mazeika2024harmbench}, and StrongREJECT \citep{souly2024strongreject} each compare automated judges against human labels of harmful compliance, and safeguard models such as Llama Guard \citep{inan2023llamaguard} and Granite Guardian \citep{padhi2024granite} are evaluated on harm taxonomies.
Refusal-string matching, introduced with the GCG attack \citep{zou2023universal}, is the classic baseline and is known to be a weak proxy.
All of these validations stay inside one domain.
Unlike this line of work, we hold the detector fixed and ask how its sensitivity changes when the misbehaviour is not harmful compliance.

\para{Bespoke instruments for subtle behaviours.}
Sycophancy is measured with dedicated prompts and answer checks \citep{perez2022discovering,sharma2023sycophancy}, and reasoning faithfulness with interventions on the chain of thought \citep{turpin2023language,lanham2023measuring}.
These protocols establish that the behaviours exist, but they do not report what general-purpose safety metrics say about the same outputs.
Building on their data---we reuse the sycophancy prompts of \citet{sharma2023sycophancy}---we run harm-style metrics on responses that these instruments would label as misaligned.

\para{Planted misalignment.}
Sleeper agents \citep{hubinger2024sleeper}, alignment-faking demonstrations \citep{greenblatt2024alignment}, and auditing games \citep{marks2025auditing} plant known misbehaviour in a model to test whether training or auditors remove or find it.
Our anchors serve a narrower and cheaper purpose: they are response-level positive controls for calibrating the per-domain sensitivity of off-the-shelf metrics, and require no model training.

\para{LLMs as judges.}
LLM judges are now a default evaluation tool \citep{zheng2023judging}, and their agreement with humans is known to vary with the prompt and the judge model.
We add a concrete case in the safety setting: the same rubric, with only its answer labels reworded, moves one judge from below chance to 95\% detection at 4\% false positives (\secref{sec:second_judge}).
